\chapter{Metodologia}
\label{chap:metodologia}

Este capítulo descreve a estrutura de software construída para implementar e comparar as seis formulações de fluxo de potência consideradas neste trabalho. A \autoref{sec:ambiente} apresenta o ambiente computacional em Julia e os pacotes utilizados. A \autoref{sec:pipeline} descreve o fluxo de execução, desde a leitura do arquivo PWF até a exportação dos resultados. A \autoref{sec:formulacoes} é o núcleo do capítulo: nela, é detalhada a formulação matemática de cada uma das seis variantes (duas baseadas em \textit{PowerModels.jl} e quatro desenvolvidas manualmente em JuMP) que, juntas, formam a progressão F0$\rightarrow$F5 usada no \autoref{chap:resultados}. Por fim, a \autoref{sec:bateria-casos} descreve o conjunto de vinte e sete casos de teste e os critérios de comparação adotados.

\section{Ambiente computacional}
\label{sec:ambiente}

Toda a implementação foi realizada na linguagem Julia~\cite{bezanson2017julia} (versão 1.x, série estável), escolhida por aliar uma sintaxe de alto nível, voltada para computação científica, a um desempenho computacional comparável ao de linguagens compiladas. O ambiente de pacotes utilizado, definido em \texttt{Project.toml} e \texttt{Manifest.toml} na raiz do repositório, agrega:

\begin{alineascomponto}
    \item \textit{PWF.jl}~\cite{pwfjl}: leitura e interpretação de arquivos no formato ANAREDE PWF, com extração de estruturas de controle habilitada via \texttt{add\_control\_data = true};
    \item \textit{PowerModels.jl}~\cite{coffrin2018powermodels}: construção da estrutura de referência (\texttt{build\_ref}) sobre os dados parseados e oferta das funções \texttt{solve\_ac\_opf} e \texttt{solve\_ac\_pf} que implementam as duas formulações de referência;
    \item \textit{JuMP.jl}~\cite{lubin2023jump}: linguagem de modelagem matemática usada para escrever as quatro formulações controladas;
    \item \textit{Ipopt.jl}~\cite{wachterBiegler2006ipopt}: interface Julia para o \textit{solver} de programação não linear baseado em método de pontos interiores;
    \item \textit{CSV.jl} e \textit{DataFrames.jl}: exportação dos resultados para arquivos CSV.
\end{alineascomponto}

A escolha do Ipopt como \textit{solver} se justifica por três fatores: (i) é \textit{open-source} e amplamente utilizado em problemas NLP de SEP; (ii) suporta nativamente expressões não lineares no JuMP via \texttt{@NLexpression} e \texttt{@NLconstraint}; e (iii) é o \textit{solver} adotado por padrão pelo \textit{PowerModels.jl} em suas formulações AC, garantindo uma comparação justa entre as referências e as formulações desenvolvidas manualmente. As opções utilizadas são \texttt{max\_iter = 3000}, \texttt{tol = 1e-5} e \texttt{print\_level = 5}. Esses valores foram escolhidos de forma permissiva por causa do porte do cenário integral do SIN (\texttt{CASO\_VER\_MAXDIU.PWF} do PAR/PEL 2027-2031, com cerca de treze mil barras), no qual tolerâncias mais apertadas comprometem a convergência.

\section{Fluxo de execução}
\label{sec:pipeline}

Cada uma das seis formulações está implementada como um \textit{script} Julia auto-contido no repositório do trabalho, seguindo a mesma estrutura sequencial:

\begin{alineascomnumero}
    \item \textbf{Leitura do PWF}: \texttt{PWF.\allowbreak parse\_file(caminho;\allowbreak\ add\_control\_data = true)} produz um dicionário com a topologia da rede, os dados de gerador, carga e ramo, além das estruturas de controle DOPC/DLIN;
    \item \textbf{Adequação da topologia}: \texttt{PowerModels.\allowbreak select\_largest\_\allowbreak component!(data)} descarta ilhas isoladas, mantendo apenas a rede principal conectada. Em seguida, se múltiplas barras forem marcadas como referência (\texttt{bus\_type == 3}), as excedentes são rebaixadas para PV (\texttt{bus\_type == 2}), preservando uma única barra \textit{slack}. Esse passo é necessário para o caso SIN, em que a leitura inicial pode marcar várias barras como referência;
    \item \textbf{Construção da estrutura de referência}: \texttt{ref = PowerModels.\allowbreak build\_ref(data)\allowbreak [:it]\allowbreak [:pm]\allowbreak [:nw]\allowbreak [0]} materializa a estrutura indexada usada para iterar sobre barras, geradores, ramos e cargas;
    \item \textbf{Construção do modelo}: criação das variáveis, expressões de fluxo e restrições no JuMP, conforme descrito na \autoref{sec:formulacoes};
    \item \textbf{Resolução}: \texttt{optimize!(model)}, com tempo de execução medido por \texttt{@elapsed};
    \item \textbf{Exportação}: dois arquivos CSV são produzidos: \texttt{resultados\_\allowbreak barras\_SIN.csv}, com $V_i$, $\theta_i$, $p^g_i$, $q^g_i$, $p^d_i$, $q^d_i$ e os \textit{slacks} de cada barra, e \texttt{resultados\_\allowbreak fluxos\_\allowbreak linhas\_SIN.csv}, com $p_{ij}, q_{ij}, p_{ji}, q_{ji}$, perdas e \textit{tap} de cada ramo.
\end{alineascomnumero}

A \autoref{fig:pipeline} resume graficamente esse fluxo sequencial.

\begin{figure}[!htb]
\Caption{\label{fig:pipeline} Fluxograma da sequência de execução das formulações}
\IBGEtab{}{
\begin{tikzpicture}[
   >=Latex, node distance=5.5mm,
   etapa/.style={draw, rounded corners, align=center, text width=9.6cm,
                 inner sep=4pt, font=\small}]
  \node[etapa] (e1) {1. Leitura do PWF com \texttt{PWF.parse\_file} (\texttt{add\_control\_data = true})};
  \node[etapa, below=of e1] (e2) {2. Adequação da topologia: maior componente conexa e barra \textit{slack} única};
  \node[etapa, below=of e2] (e3) {3. Construção da estrutura de referência (\texttt{build\_ref})};
  \node[etapa, below=of e3] (e4) {4. Construção do modelo: variáveis, expressões de fluxo e restrições no JuMP};
  \node[etapa, below=of e4] (e5) {5. Resolução com \texttt{optimize!} e medição do tempo de execução};
  \node[etapa, below=of e5] (e6) {6. Exportação dos resultados em dois arquivos CSV (barras e ramos)};
  \draw[->] (e1) -- (e2);
  \draw[->] (e2) -- (e3);
  \draw[->] (e3) -- (e4);
  \draw[->] (e4) -- (e5);
  \draw[->] (e5) -- (e6);
\end{tikzpicture}
}{
\Fonte{elaborado pelo autor.}
}
\end{figure}

Os elos HVDC, presentes em vários casos de teste, têm um tratamento particular nas formulações baseadas em \textit{PowerModels} e nas desenvolvidas manualmente: as injeções $(p^{dc}_{ft}, q^{dc}_{ft})$ em cada extremidade são lidas do PWF e \textbf{fixadas} via \texttt{JuMP.\allowbreak fix(...;\allowbreak\ force=true)}. Essa decisão segue a recomendação do CEPEL de tratar elos HVDC como condições de contorno do FP AC: as injeções já vêm pré-especificadas pelo fluxo base e liberá-las apenas amplia o Jacobiano sem trazer informação nova. O retificador recebe $p^g$ negativo e o inversor $p^g$ positivo.

\section{Formulações implementadas}
\label{sec:formulacoes}

A progressão de seis formulações é identificada pelas siglas da \autoref{tab:progressao-formulacoes}, que resumem os controles incorporados em cada variante.

\begin{table}[!h]
\captionsetup{width=15cm}
\Caption{\label{tab:progressao-formulacoes} Progressão das formulações implementadas}
\IBGEtab{}{
    \begin{tabular}{cllll}
        \toprule
        Sigla & Arquivo & Motor & Tipo & Acrescenta \\
        \midrule \midrule
        F0 & \texttt{(F0)OPF\_PM.jl} & PowerModels & FPO AC & referência de FPO \\
        F1 & \texttt{(F1)PF\_PM.jl} & PowerModels & FP AC & referência de FP \\
        F2 & \texttt{(F2)QLIM+VLIM.jl} & JuMP (manual) & FP controlado & QLIM + VLIM \\
        F3 & \texttt{(F3)QLIM+VLIM+CSCA.jl} & JuMP (manual) & FP controlado & + CSCA \\
        F4 & \texttt{(F4)QLIM+VLIM+CSCA+CTAP.jl} & JuMP (manual) & FP controlado & + CTAP \\
        F5 & \texttt{(F5)+DERA.jl} & JuMP (manual) & FP controlado & + DERA \\
        \bottomrule
    \end{tabular}
}{
    \Fonte{elaborado pelo autor.}
}
\end{table}

A progressão é estritamente \emph{aditiva}: cada arquivo numerado parte do anterior e acrescenta um único bloco de variáveis e restrições. A leitura no \autoref{chap:resultados} fica simplificada porque qualquer divergência entre dois \textit{scripts} consecutivos pode ser atribuída exclusivamente ao novo controle introduzido.

\subsection{Formulação F0: FPO AC de referência}
\label{subsec:f14}

A formulação F0 é uma chamada direta a \texttt{PowerModels.\allowbreak solve\_ac\_opf(data,\allowbreak\ optimizer)}, sem reconstrução manual de variáveis ou restrições. O modelo subjacente é exatamente o FPO AC canônico descrito na Equação \eqref{eq:fpo}: minimiza $\sum_k c_k(p^g_k)$ sujeito ao balanço nodal e aos limites operacionais em $V$, $p^g$, $q^g$ e $|s_{ij}|$. Esta formulação serve como referência de FPO: não incorpora nenhuma das ações de controle ANAREDE da \autoref{sec:acoes-controle}, mas oferece o ponto de operação economicamente ótimo da rede tal como definido pela literatura clássica.

Cabe aqui uma ressalva metodológica importante. Não estão disponíveis dados reais sobre os custos de operação das fontes do SIN, e os arquivos PWF utilizados não trazem curvas de custo individuais por gerador. Na ausência dessa informação, os coeficientes de custo atribuídos aos geradores são idênticos entre si e, como a demanda total é fixa, minimizar o custo total de geração torna-se equivalente a minimizar a geração total --- ou seja, a minimizar as perdas ativas da rede. A saída desta formulação deve, portanto, ser interpretada estritamente como um ponto de operação de mínimas perdas ativas, não sendo adequada para análises de viabilidade econômica neste contexto.

\subsection{Formulação F1: FP AC clássico}
\label{subsec:f15}

A formulação F1 é uma chamada a \texttt{PowerModels.\allowbreak solve\_ac\_pf(data,\allowbreak\ optimizer)}. Diferentemente de F0, não há otimização: a função objetivo é vazia e o problema reduz-se à resolução do sistema algébrico de FP do tipo Newton, com $V$ e $\theta$ como incógnitas em barras PQ, $\theta$ em barras PV e nada em barras \textit{slack}. F1 representa o ponto de operação ``puro'' da rede, sem qualquer otimização de custos ou ação de controle.

\subsection{Formulação F2: FP controlado QLIM + VLIM}
\label{subsec:f16}

A partir de F2, o trabalho abandona o \textit{PowerModels.jl} como motor de resolução e passa a construir o modelo diretamente em JuMP. É importante mencionar que a leitura dos dados continua sendo feita pelo \textit{parser} do \textit{PWF.jl}, e a estrutura de referência indexada do \textit{PowerModels.jl} (\texttt{build\_ref}) continua sendo utilizada para iterar sobre barras, geradores, ramos e cargas; o que muda é exclusivamente a montagem das variáveis, restrições e função objetivo, agora feita manualmente. As variáveis primárias são as mesmas do FP clássico:

\begin{equation}
\label{eq:vars-base}
V_i \in [V^{\min}_i, V^{\max}_i], \quad \theta_i \in \mathbb{R}, \quad p^g_k \in [p^{g,\min}_k, p^{g,\max}_k], \quad q^g_k \in [q^{g,\min}_k, q^{g,\max}_k],
\end{equation}

\noindent inicializadas com os valores do PWF (\texttt{start =\allowbreak\ ref[:bus][i]\allowbreak ["vm"]}, etc.). Os fluxos $p_{ij}, q_{ij}$ são definidos como \texttt{@NLexpression} sobre as Equações \eqref{eq:p-ramo}--\eqref{eq:q-ramo}, e as Equações \eqref{eq:balanco-p}--\eqref{eq:balanco-q} são impostas via \texttt{@NLconstraint}.

A novidade de F2 está na introdução das variáveis de folga:

\begin{equation}
\label{eq:slacks-vlim-qlim}
s^v_i \in \mathbb{R}, \;\; \forall i \in \mathcal{N}^{gen}, \qquad s^q_l \in \mathbb{R}, \;\; \forall l \in \mathcal{L},
\end{equation}

\noindent e da restrição de \textit{setpoint} de tensão para cada barra de geração, que implementa o QLIM: com $q^g_k$ mantido rigidamente dentro de $[q^{g,\min}_k, q^{g,\max}_k]$, a folga $s^v_i$ absorve o desvio de tensão quando o limite de reativo impede sustentar o \textit{setpoint}:

\begin{equation}
\label{eq:setpoint-vlim}
V_i = V^{ref}_i + s^v_i, \qquad \forall i \in \mathcal{N}^{gen},
\end{equation}

\noindent onde $V^{ref}_i$ é o valor de tensão lido do PWF. A folga de potência reativa $s^q_l$, que implementa o VLIM, entra diretamente no balanço reativo da Equação \eqref{eq:balanco-q}, que em F2 passa a ser escrita como:

\begin{equation}
\label{eq:balanco-q-folga}
\sum_{(i,j) \in \mathcal{B}_i} q_{ij} \;=\; \sum_{k \in \mathcal{G}_i} q^g_k \;-\; \sum_{l \in \mathcal{L}_i} \big( q^d_l + s^q_l \big) \;+\; b^{sh}_i V_i^2, \qquad \forall i \in \mathcal{N},
\end{equation}

\noindent isto é, $s^q_l$ permite que a demanda reativa efetiva da carga $l$ se desvie do valor nominal apenas o necessário para que a tensão da barra permaneça dentro dos limites. Essa é a mesma formulação de VLIM adotada no \textit{ControlPowerFlow.jl} \cite{chavarry2024}, com uma diferença de escrita: naquele pacote, a demanda reativa é promovida a variável de decisão explícita, amarrada ao valor especificado pela restrição $q^d_l = q^{d,spec}_l + s^q_l$; aqui, essa variável intermediária é eliminada por substituição direta no balanço nodal, resultando em um modelo matematicamente equivalente com menos variáveis. Em barras \textit{slack}, $\theta_i$ é fixado em zero via \texttt{fix(va[i],\allowbreak\ 0.0;\allowbreak\ force=true)}. Em geradores \textit{slack}, os limites de $p^g_k$ são removidos para acomodar o balanço global de potência ativa.

A função objetivo é puramente de penalidade quadrática:

\begin{equation}
\label{eq:obj-16}
\min \;\; \rho \, \sum_{i \in \mathcal{N}^{gen}} (s^v_i)^2 \;+\; \rho \, \sum_{l \in \mathcal{L}} (s^q_l)^2,
\end{equation}

\noindent com $\rho = 10^6$. Esse valor de penalidade é alto o suficiente para que, sempre que o problema for factível com $s^v = s^q = 0$, a solução ótima respeite os \textit{setpoints}, mas baixo o suficiente para não comprometer o condicionamento numérico do problema.

\subsection{Formulação F3: adiciona CSCA}
\label{subsec:f17}

A formulação F3 parte de F2 e introduz o controle de bancos \textit{shunt} chaveados. Para cada banco $s \in \mathcal{S}$, é criada a variável contínua

\begin{equation}
\label{eq:csca-var}
b^{sh}_s \in [b^{sh,\min}_s, b^{sh,\max}_s],
\end{equation}

\noindent e o termo $b^{sh}_s V_i^2$ é adicionado ao balanço reativo da Equação \eqref{eq:balanco-q} na barra $i$ hospedeira. Diferentemente de F2, em que o desvio de tensão é tratado por uma única folga de igualdade $s^v_i$ aplicada ao \textit{setpoint}, F3 reorganiza o tratamento de tensão das barras controladas em dois regimes: nas barras cujo intervalo de controle é praticamente um ponto ($|V^{\max}_i - V^{\min}_i| < 10^{-5}$) mantém-se a restrição de igualdade $V_i = V^{ref}_i + s^v_i$; nas demais, a tensão é mantida dentro da faixa $[V^{\min}_i, V^{\max}_i]$ por meio de um par de folgas \emph{unilaterais} não negativas, $s^{v,\downarrow}_i$ no limite inferior e $s^{v,\uparrow}_i$ no superior:

\begin{equation}
\label{eq:csca-slack-v}
V_i \ge V^{\min}_i - s^{v,\downarrow}_i, \qquad V_i \le V^{\max}_i + s^{v,\uparrow}_i, \qquad s^{v,\downarrow}_i, s^{v,\uparrow}_i \ge 0.
\end{equation}

\noindent Além disso, é introduzida uma folga $s^b_s$ que liga $b^{sh}_s$ ao seu valor nominal lido do PWF ($b^{sh}_s = b^{sh,\text{nom}}_s + s^b_s$), penalizada com peso menor. A função objetivo de F3 passa, portanto, a ter quatro blocos:

\begin{equation}
\label{eq:obj-17}
\min \;\; \rho \sum_{i \in \mathcal{N}^{ctrl}} \big[(s^v_i)^2 + (s^{v,\uparrow}_i)^2 + (s^{v,\downarrow}_i)^2 \big] \;+\; \rho \sum_{l \in \mathcal{L}} (s^q_l)^2 \;+\; \rho_b \sum_{s \in \mathcal{S}} (s^b_s)^2,
\end{equation}

\noindent com $\rho = 10^6$ (folgas de tensão e de reativo) e $\rho_b = 10^4$ (folga de \textit{shunt}). O peso menor sobre $s^b_s$ é deliberado: garante que o \textit{solver} prefira ajustar o banco \textit{shunt} dentro de sua faixa a violar os \textit{setpoints} de tensão. O efeito é que o \textit{solver} ajusta o valor de cada banco \textit{shunt} para minimizar o desvio das tensões em relação aos \textit{setpoints}, exatamente o papel que o CSCA do ANAREDE desempenha por meio de chaveamentos discretos. O caráter discreto do controle real é deixado para uma etapa de discretização posterior, fora do escopo deste trabalho.

\subsection{Formulação F4: adiciona CTAP}
\label{subsec:f18}

A formulação F4 acrescenta a F3 o controle automático de \textit{tap} de transformadores. Para cada ramo $(i,j) \in \mathcal{B}^{tap}$ classificado no PWF como transformador com OLTC, a razão $t_{ij}$, antes constante, passa a ser variável:

\begin{equation}
\label{eq:ctap-var}
t_{ij} \in [t^{\min}_{ij}, t^{\max}_{ij}],
\end{equation}

\noindent e as expressões de fluxo \eqref{eq:p-ramo}--\eqref{eq:q-ramo} são reescritas com $t_{ij}$ no denominador como variável, o que adiciona um fator não linear adicional às equações já não lineares do balanço. F4 não introduz novas folgas nem novos termos na função objetivo: a regulação de tensão associada ao \textit{tap} é absorvida pelas mesmas folgas de tensão $s^v_i$, $s^{v,\uparrow}_i$ e $s^{v,\downarrow}_i$ já presentes em F3 (Equação \eqref{eq:obj-17}), agora também sensíveis ao novo grau de liberdade $t_{ij}$. O resultado é que o \textit{solver} encontra simultaneamente os \textit{taps} e os pontos de operação compatíveis, em uma única resolução NLP. Isso contrasta com o ANAREDE, que itera externamente entre o FP e o ajuste dos \textit{taps}.

\subsection{Formulação F5: adiciona DERA}
\label{subsec:f19}

A formulação F5 acrescenta a F4 o esquema hierárquico de corte/alívio de carga descrito na \autoref{subsec:dera}. Para cada carga $l \in \mathcal{L}$ são introduzidas duas variáveis de decisão:

\begin{equation}
\label{eq:dera-vars}
\Delta p^d_l \in [0, p^d_l], \qquad \Delta q^d_l \in \mathbb{R},
\end{equation}

\noindent representando, respectivamente, o corte de potência ativa e o corte de potência reativa daquela carga. As duas variáveis ficam acopladas por uma restrição de fator de potência constante:

\begin{equation}
\label{eq:dera-fp}
\Delta q^d_l \;=\; \Delta p^d_l \, \frac{q^d_l}{p^d_l}, \qquad \forall l \in \mathcal{L} \text{ com } p^d_l > 10^{-4}.
\end{equation}

\noindent Cargas com $p^d_l \le 10^{-4}$ (efetivamente nulas) têm $\Delta p^d_l$ e $\Delta q^d_l$ fixados em zero via \texttt{fix(...;\allowbreak\ force=true)}, eliminando-as do problema. Os cortes entram nos balanços nodais ativo e reativo \eqref{eq:balanco-p}--\eqref{eq:balanco-q} subtraindo das demandas nominais: $p^d_l \rightarrow p^d_l - \Delta p^d_l$ e $q^d_l \rightarrow q^d_l - \Delta q^d_l$.

A novidade de F5 na função objetivo é um termo \emph{linear} ponderado pelo nível de tensão da barra hospedeira:

\begin{equation}
\label{eq:obj-19}
\min \;\; \underbrace{\rho \!\! \sum_{i \in \mathcal{N}^{ctrl}} \!\! \big[(s^v_i)^2 + (s^{v,\uparrow}_i)^2 + (s^{v,\downarrow}_i)^2 \big] + \rho \sum_{l \in \mathcal{L}} (s^q_l)^2 + \rho_b \sum_{s \in \mathcal{S}} (s^b_s)^2}_{\text{termos herdados de F2--F4}} \;+\; \sum_{l \in \mathcal{L}} w_l \, \Delta p^d_l,
\end{equation}

\noindent com pesos $w_l$ definidos por faixa de tensão base da barra hospedeira de $l$:

\begin{equation}
\label{eq:dera-pesos}
w_l \;=\; \begin{cases}
10^3, & V^{base}_l < 69\,\text{kV} \quad (\text{subtransmissão e distribuição}), \\
5 \times 10^3, & 69 \le V^{base}_l \le 230\,\text{kV} \quad (\text{transmissão}), \\
10^4, & V^{base}_l > 230\,\text{kV} \quad (\text{extra alta tensão, EAT}).
\end{cases}
\end{equation}

\noindent Em F5, o coeficiente de penalidade dos \textit{slacks} de tensão é reduzido de $\rho = 10^6$ (valor usado em F2--F4) para $\rho = 10^5$, e $\rho_b = 10^3$, de modo que o termo linear de corte de carga seja efetivamente comparável aos termos quadráticos de \textit{slack}. Com a penalidade original a um milhão, os \textit{slacks} de tensão dominariam a função objetivo e o corte de carga praticamente nunca seria acionado. O efeito esperado da formulação é o seguinte: enquanto o \textit{solver} consegue acomodar todas as restrições de tensão e \textit{setpoints} apenas com os \textit{slacks} $s^v, s^q, s^b$, o corte $\Delta p^d_l$ permanece em zero e F5 reproduz F4. Quando uma carga estaria sendo servida em uma barra cuja tensão violaria intoleravelmente o \textit{setpoint}, o \textit{solver} passa a preferir um corte parcial dessa carga (começando pelas barras de tensão mais baixa, onde $w_l = 10^3$) em vez de aumentar as folgas de tensão (que custam $\rho = 10^5$ ao quadrado). O esquema reproduz, no nível de função objetivo, a hierarquia de prioridade de corte de carga adotada em emergência pelo operador do sistema.

\subsection{Formulação auxiliar FDC: fluxo de potência DC do \textit{PowerModels}}
\label{subsec:fdc}

Além da progressão principal F0$\rightarrow$F5, foi implementado um \textit{script} auxiliar (\texttt{(FDC).jl}) que resolve o caso integral do SIN sob a \emph{aproximação DC} do \textit{PowerModels} (\texttt{solve\_dc\_pf}), descrita na \autoref{subsec:relaxacoes-fpo}. A motivação dessa formulação isolada é responder a uma pergunta deixada em aberto pela discussão de convergência do caso SIN no \autoref{chap:resultados}: dado que nenhuma das seis formulações AC converge sobre \texttt{CASO\_\allowbreak VER\_MAXDIU.PWF} dentro do limite de iterações, a versão linearizada do mesmo problema é factível?

A FDC herda integralmente o fluxo de execução de F0/F1, as entradas e as saídas são as mesmas (PWF, base 100\,MVA), e zera os fluxos reativos e perdas ativas (desprezados por definição). A FDC não substitui nenhuma das seis formulações principais e não é executada sobre o conjunto completo de casos de teste: trata-se de uma verificação isolada sobre o caso SIN, cujos resultados são discutidos na \autoref{subsec:fdc-resultado}.

\section{Conjunto de casos de teste e critérios de comparação}
\label{sec:bateria-casos}

\subsection{Casos de teste}
\label{subsec:casos-teste}

O conjunto de casos de teste utilizado nas comparações do \autoref{chap:resultados} é composto por vinte e sete arquivos PWF, agrupados em cinco faixas de porte: redes pedagógicas, redes de médio porte, um cenário integral do SIN, recortes reduzidos do SIN e casos sintéticos de estresse. A escolha foi pensada para varrer dois eixos simultaneamente: o porte da rede (de três a 13~338 barras) e a presença explícita de cada uma das ações de controle modeladas neste trabalho.

A convenção de nomenclatura dos arquivos pedagógicos segue o padrão \texttt{<base>\allowbreak\_<sufixo>}, em que \texttt{<base>} indica o número de barras e a família (\texttt{3bus} para a rede triangular clássica de três barras, \texttt{3busfrank}/\allowbreak \texttt{4busfrank}/\allowbreak \texttt{5busfrank} para variantes da família \textit{frank} distribuídas junto ao próprio \textit{PWF.jl} como casos didáticos, \texttt{9bus} para a rede de nove barras de Anderson-Fouad) e \texttt{<sufixo>} indica precisamente o elemento ou controle que foi adicionado à versão base do arquivo. Os sufixos seguem o nome da seção PWF ou da palavra-chave da opção de controle introduzida, conforme detalhado na \autoref{tab:bateria-casos}.

\begin{table}[!htb]
\captionsetup{width=16cm}
\Caption{\label{tab:bateria-casos} Conjunto de casos de teste: porte e diferença em relação à versão base}
\IBGEtab{}{
    \footnotesize
    \setlength{\tabcolsep}{4pt}
    \renewcommand{\arraystretch}{1.1}
    \begin{tabular}{p{4.4cm}rp{9.6cm}}
        \toprule
        Arquivo & Barras & Diferença em relação à base / observação \\
        \midrule \midrule
        \multicolumn{3}{l}{\textit{Família 3bus} (rede triangular clássica de três barras)} \\
        \midrule
        \texttt{3bus.pwf}                & 3      & versão base, apenas seções DBAR, DLIN, DCTE \\
        \texttt{3bus\_DBSH.pwf}          & 3      & + seção DBSH (banco \textit{shunt} chaveável detalhado) \\
        \texttt{3bus\_DCER.pwf}          & 3      & + seção DCER (compensador estático de reativos, SVC) \\
        \texttt{3bus\_DCSC.pwf}          & 3      & + seção DCSC (capacitor série controlável, CSC) \\
        \texttt{3bus\_DSHL.pwf}          & 3      & + seção DSHL (\textit{shunt} de linha, no terminal de envio ou recepção) \\
        \texttt{3bus\_DCline.pwf}        & 3      & + elo HVDC completo (seções DELO, DCBA, DCLI, DCNV, DCCV) \\
        \texttt{3bus\_shunt\_fields.pwf} & 3      & combina simultaneamente DBSH e DCER, exercendo o parser em campos \textit{shunt} concorrentes \\
        \texttt{3bus\_corrections.pwf}   & 3      & variante com campos numéricos modificados para testar correções de parser; dispara \textbf{ERR} \\
        \midrule
        \multicolumn{3}{l}{\textit{Família \textit{frank}} (variantes do livro-texto de Frank)} \\
        \midrule
        \texttt{3busfrank.pwf}                     & 3 & versão base sem controles \\
        \texttt{3busfrank\_continuous\_shunt.pwf}  & 3 & + DBSH com passo contínuo (relaxa a discretização do CSCA) \\
        \texttt{3busfrank\_qlim.pwf}               & 3 & ativa \textbf{QLIM} na seção DOPC (controle de limite de reativo) \\
        \texttt{4busfrank\_vlim.pwf}               & 4 & ativa \textbf{VLIM} na seção DOPC (controle de magnitude de tensão) \\
        \texttt{5busfrank.pwf}                     & 5 & versão base sem controles \\
        \texttt{5busfrank\_csca.pwf}               & 5 & ativa \textbf{CSCA} na seção DOPC (\textit{shunt} chaveável automático) \\
        \texttt{5busfrank\_ctap.pwf}               & 5 & ativa \textbf{CTAP} na seção DOPC (\textit{tap} de transformador) \\
        \texttt{5busfrank\_ctaf.pwf}               & 5 & ativa \textbf{CTAF} na seção DOPC (\textit{tap} de transformador defasador, variante do CTAP) \\
        \texttt{5busfrank\_cphs.pwf}               & 5 & ativa \textbf{CPHS} na seção DOPC (controle de fase) \\
        \midrule
        \multicolumn{3}{l}{\textit{Família 9bus} (Anderson-Fouad)} \\
        \midrule
        \texttt{9bus.pwf}                     & 9 & versão base \\
        \texttt{9bus\_transformer\_fields.pwf}& 9 & + seções DCTR, DGLT e DTPF (campos detalhados de transformador e grupos limite de tensão) \\
        \midrule
        \multicolumn{3}{l}{\textit{Redes de médio porte}} \\
        \midrule
        \texttt{300bus.pwf}              & 300    & IEEE 300-bus, com múltiplas barras de geração, elos HVDC e bancos \textit{shunt} \\
        \texttt{500bus.pwf}              & 500    & rede sintética de 500 barras com SVCs e múltiplos OLTC \\
        \midrule
        \multicolumn{3}{l}{\textit{Cenário integral do SIN}} \\
        \midrule
        \texttt{CASO\_VER\_MAXDIU.PWF}   & 13\,338 & SIN integral, Verão 2027/2028 Máxima Diurna do PAR/PEL 2027-2031 do ONS \\
        \midrule
        \multicolumn{3}{l}{\textit{Recortes reduzidos do SIN (derivados de \texttt{CASO\_VER\_MAXDIU})}} \\
        \midrule
        \texttt{caso\_red.pwf}           & 2\,599 & recorte do subsistema Nordeste completo, com 15 barras de fronteira no eixo PI-BA (\autoref{subsec:reducao-fronteira}) \\
        \texttt{caso\_red2.pwf}          & 1\,033 & recorte ainda mais reduzido (CE + RN + PB) do subsistema acima \\
        \midrule
        \multicolumn{3}{l}{\textit{Casos sintéticos de estresse do \textit{parser} e dos controles}} \\
        \midrule
        \texttt{test\_defaults.pwf}      & 3 & exercita campos opcionais não preenchidos (\textit{defaults}) em quase todas as seções (DBAR, DBSH, DCBA, DELO, etc.) \\
        \texttt{test\_line\_shunt.pwf}   & 3 & exercita combinações de DBSH e DSHL na mesma linha \\
        \texttt{test\_system.pwf}        & 3 & cenário sintético mínimo (apenas DBAR e DLIN) usado como verificação de consistência \\
        \bottomrule
    \end{tabular}
}{
    \Fonte{elaborado pelo autor. ``+'' indica seção adicionada em relação ao arquivo base imediatamente acima da linha; \textbf{negrito} em letras maiúsculas indica palavra-chave ativada na linha de opções do registro DOPC do PWF.}
}
\end{table}

Para cada caso, todas as seis formulações são executadas com a mesma configuração de \textit{solver}, a mesma versão do parser \textit{PWF.jl} e a mesma máquina, garantindo a comparabilidade dos tempos e dos status de convergência.

\subsection{Recortes reduzidos do SIN: método de fronteira híbrida}
\label{subsec:reducao-fronteira}

Os arquivos \texttt{caso\_red.pwf} e \texttt{caso\_red2.pwf} foram obtidos de uma colaboração externa (M. Sc.\ Isabela Metzker, ONS) a partir do caso \texttt{CASO\_VER\_MAXDIU.PWF}, que reproduz a condição operativa Verão 2027/2028 Máxima Diurna do ciclo PAR/PEL 2027-2031 do ONS. O \texttt{caso\_red} preserva o subsistema Nordeste completo (2 599 barras) e o \texttt{caso\_red2} restringe-se a Ceará, Rio Grande do Norte e Paraíba (1 033 barras), agrupamento natural justificado pela concentração de geração eólica em CE/RN e pela contiguidade elétrica entre RN e PB.

A redução adota um esquema híbrido de fronteira, sem ser uma equivalência tipo Ward formal:

\begin{alineascomnumero}
    \item para cada barra de fronteira, a carga e a geração externas que originalmente fluíam pela barra são compactadas em valores equivalentes embutidos diretamente nos campos $p^d, q^d, p^g, q^g$ da seção DBAR. Em particular, as quinze barras de fronteira do \texttt{caso\_red}, localizadas no eixo PI-BA do SIN, recebem cargas positivas (sumidouros NE $\rightarrow$ SE: \texttt{R.EGUA}, \texttt{RE-LZ-CAP}, \texttt{POCIII}, \texttt{M.NETO}, \texttt{IGAPOR}, \texttt{BJLAP2}) ou cargas negativas (fontes N $\rightarrow$ NE: \texttt{TERES2}, \texttt{B.ESPE}, \texttt{RG-CO2}, \texttt{RG-CO1}), representando a injeção líquida de cada caminho externo cortado;
    \item para cada par de barras de fronteira que originalmente se conectava pelo lado externo, é inserido um ramo direto na seção DLIN com impedância equivalente ao caminho cortado. No \texttt{caso\_red}, vinte e um ramos físicos foram removidos e cinquenta e seis ramos equivalentes foram inseridos.
\end{alineascomnumero}

A consequência direta para a análise de resultados é que as perdas ativas dos casos reduzidos incluem a dissipação real nessas impedâncias equivalentes; isto é, parte do que aparece como ``perda'' corresponde à dissipação de potência ativa associada às impedâncias equivalentes resultantes da redução de fronteira do sistema, e não a perdas internas do próprio recorte. As seções DCSC, DBSH, DCER, DCTR, DGER, DGLT e DGBT do PWF original foram preservadas; as opções de controle CTAP e CPHS foram explicitamente desligadas pela equipe que produziu o recorte (DOPC com \texttt{CTAP D CPHS D}), enquanto QLIM, VLIM, CREM e AREG permanecem ativadas. As seções HVDC do SIN original (DELO, DCBA, DCLI, DCNV, DCCV) não aparecem nos recortes, uma vez que os elos Itaipu, Madeira e Belo Monte não estão na região Nordeste.

\subsection{Critérios de comparação}
\label{subsec:criterios}

Para cada par (caso, formulação), três informações são registradas:

\begin{alineascomnumero}
    \item \textbf{Status de convergência}: codificado em oito categorias, a saber, \textbf{OK} (LOCALLY\_SOLVED), \textbf{INF} (INFEASIBLE/LOCALLY\_INFEASIBLE), \textbf{ITL} (ITERATION\_LIMIT), \textbf{TIM} (TIME\_LIMIT), \textbf{KIL} (TIMEOUT\_KILLED), \textbf{CRA} (SUBPROCESS\_CRASH), \textbf{INV} (INVALID\_MODEL) e \textbf{ERR} (OTHER\_ERROR);
    \item \textbf{\textit{Cluster} de equivalência numérica}: para cada caso, as formulações que convergiram são particionadas em grupos onde duas formulações estão no mesmo grupo se, para todas as barras, $|\Delta V| \le 10^{-4}$ e $|\Delta \theta| \le 10^{-4}$, com equivalência análoga em $p^g, q^g$ e nos fluxos de ramos. \textit{Clusters} distintos representam soluções numericamente diferentes da mesma rede;
    \item \textbf{Análise barra a barra}: para os casos em que pelo menos uma das referências F0 ou F1 convergiu, são identificadas as barras nas quais as formulações F2 a F5 divergem de \emph{ambas} as referências por mais de $10^{-4}$ em módulo, em $V$ ou $\theta$. Essa análise mostra \emph{onde} as ações de controle deslocam o ponto de operação em relação à referência sem controles.
\end{alineascomnumero}

A execução do conjunto de casos de teste, a agregação dos resultados e a geração dos relatórios de convergência e de divergência barra a barra são automatizadas por \textit{scripts} auxiliares do repositório do trabalho. As tabelas e análises do \autoref{chap:resultados} são extraídas do relatório consolidado produzido por esse fluxo de processamento.
